\subsection{Descompresión selectiva y estrategia de acceso}

En una arquitectura densa tradicional, cualquier compuerta opera directamente sobre el arreglo completo de amplitudes. En contraste, cuando el
estado está comprimido por bloques, la política de acceso debe asegurar que solo se materialicen en memoria densa los bloques realmente tocados por
la operación actual. De este modo, una compuerta de uno o dos qubits no desencadena una descompresión global, sino un proceso local de adquisición
de bloques, modificación y posterior recompresión.

Esta idea conduce a una estrategia de descompresión selectiva. Cada acceso a una amplitud o a un conjunto de amplitudes se traduce en la identificación
de los bloques involucrados. Si un bloque ya se encuentra disponible en forma descomprimida, la operación se realiza directamente sobre ese buffer.
Si el bloque aún está almacenado en forma comprimida, primero se descomprime, se ejecuta la actualización necesaria y luego se conserva temporalmente
en memoria de trabajo para evitar recomprimirlo de inmediato si es probable que vuelva a usarse en operaciones posteriores.

Desde el punto de vista algorítmico, esta estrategia es especialmente relevante para compuertas que inducen accesos repetidos a subconjuntos del
estado, así como para secuencias cortas de compuertas sobre qubits cercanos. En esos casos, mantener un subconjunto pequeño del vector en memoria
descomprimida puede amortiguar el costo de las transiciones entre representación comprimida y representación densa.
